\section{The first primitive-collision transition}

Write
\[
 S_{q,L}=\{m q^{-1}\bmod h_L:m\in\Mset_L(q)\}.
\]
Since $L\le4$, the cutoff lies between $L$ and $2L-1$, so every allowed multiplier
has magnitude at most seven.  The stable assumption $Q\ge8$ therefore gives
$|m|<q$ for every active prime.

\begin{theorem}[Primitive transition]
\label{thm:transition}
For distinct primes $q_1,q_2\in\Qshell$, the supports $S_{q_1,L}$ and
$S_{q_2,L}$ are disjoint for $L=1,2,3$.  For $L=4$, every collision is, up to
exchanging the primes and changing both signs,
\begin{equation}
 7p+3r=16Q,\qquad m_p=3,\qquad m_r=-7.
 \label{eq:resonance}
\end{equation}
Conversely, a prime pair satisfying \eqref{eq:resonance}, the literal cutoff
conditions, and $(21,16Q)=1$ produces exactly two shared coordinates.
\end{theorem}

\begin{proof}
A shared coordinate is equivalent to
\[
 m_1q_2-m_2q_1\equiv0\pmod{4LQ}.
\]
Equal signs give a difference of magnitude less than $(4L-3)Q<4LQ$.  Zero
difference would force one active prime to divide a nonzero multiplier of magnitude at
most seven.  Thus the signs are opposite, and a collision reduces to
\begin{equation}
 a q_2+b q_1=4LQ,
 \qquad a,b>0.
 \label{eq:positive-collision}
\end{equation}
Primitivity makes $a$ and $b$ odd, while the prime-shell inequalities imply
$2L<a+b<4L$.

For $L=1$ no pair remains.  At $L=2$, only $(3,3)$ remains, but both cutoff
conditions force $q_i\ge3Q/2$, making the left side of
\eqref{eq:positive-collision} at least $9Q>8Q$.  At $L=3$, multiplier $3$ is
nonprimitive modulo $12Q$; the only remaining pair is $(5,5)$, whose cutoffs force
at least $50Q/3>12Q$.

At $L=4$, the possible positive primitive multipliers are $1,3,5,7$.  The size and
cutoff constraints reduce the table to $(3,7)$, $(5,5)$, $(5,7)$, and $(7,7)$.
The last two exceed $16Q$ after their own cutoff bounds are imposed.  The pair
$(5,5)$ would imply $5(q_1+q_2)=16Q$, but using multiplier five primitively requires
$5\nmid Q$.  Only $(3,7)$ and its exchange survive.  The two global sign choices give
the two coordinates, which cannot coincide because $16Q\nmid6$.
\end{proof}

\begin{table}[t]
\centering
\small
\begin{tabular}{ccl}
\toprule
$L$ & primitive collision? & first obstruction or surviving type \\
\midrule
1 & no & cutoff-one block orthogonality \\
2 & no & $(3,3)$ violates its cutoff lower bound \\
3 & no & $m=4$ is nonprimitive; $(5,5)$ is too large \\
4 & yes & only $7p+3r=16Q$ with $(3,-7)$ \\
\bottomrule
\end{tabular}
\caption{The first nontrivial cutoff is not the first primitive collision.  The
integer-dilation transition occurs at $L=4$.}
\label{tab:transition}
\end{table}

At $Q=25$, the active pair $(p,r)=(37,47)$ satisfies
$7\cdot37+3\cdot47=400$.  The inverses are $37^{-1}=173$ and
$47^{-1}=383$ modulo $400$, so the shared coordinates are
\[
 3\cdot173\equiv-7\cdot383\equiv119,
 \qquad
 -3\cdot173\equiv7\cdot383\equiv281\pmod{400}.
\]
This is the first stable collision found by the exact boundary census, and the theorem
shows that its type is the only possible one.
